\section{Introduction}
\label{sec:introduction}

For an integer $n\ge2$, write $[n]=\{1,\ldots,n\}$ and let $E_n=\binom{[n]}2$ be the edge set of the complete simple graph on $n$ vertices. We identify a graph with an element of $\{0,1\}^{E_n}$. A \emph{graph property} is a family $\mathcal P\subseteq\{0,1\}^{E_n}$ invariant under every permutation of $[n]$. It is \emph{increasing}, or \emph{monotone}, if adding edges preserves membership, and \emph{nontrivial} if it is neither empty nor the whole space.

\noindent\textit{Remark (notation).} $\mathcal P\subseteq\{0,1\}^{E_n}$ is a graph family invariant under every permutation of $[n]$. In Theorem~\ref{thm:threshold}: $\mathcal P$ is increasing under addition of edges and satisfies $\varnothing\ne\mathcal P\ne\{0,1\}^{E_n}$.


Write $\mu_p(\mathcal P)$ for the probability that a graph whose edges are present independently with probability $p$ belongs to $\mathcal P$. For $0<a<1$, define
\[
 p_a(\mathcal P)=\inf\{p\in[0,1]:\mu_p(\mathcal P)\ge a\}.
\]
All logarithms in this paper are natural. In particular, $n$ always counts vertices, not edge coordinates.

\noindent\textit{Remark (notation).} For $0<a<1$, $p_a(\mathcal P)=\inf\{p\in[0,1]:\mu_p(\mathcal P)\ge a\}$ whenever this defining set is nonempty. In Theorem~\ref{thm:threshold}: $p_a(\mathcal P)\in(0,1)$ is the unique solution of $\mu_{p_a(\mathcal P)}(\mathcal P)=a$ for every $0<a<1$.


\begin{theorem}[Friedgut--Kalai sharp-threshold conjecture]
\label{thm:threshold}
For every integer $n\ge2$, every nontrivial increasing family $\mathcal P$ of graphs on $n$ vertices invariant under all vertex permutations, and every $0<\varepsilon<1/2$,
\begin{equation}
\label{eq:threshold}
 p_{1-\varepsilon}(\mathcal P)-p_\varepsilon(\mathcal P)
 \le \frac{2^{19}}{(\log n)^2}\log\frac1{2\varepsilon}.
\end{equation}
\end{theorem}

The influence theorem of Kahn, Kalai, and Linial~\cite{KKL88}
connected Fourier analysis with the sensitivity of Boolean functions
to individual coordinates. Bourgain, Kahn, Kalai, Katznelson, and
Linial extended the influence theorem to product
spaces~\cite{BKKKL92}. Friedgut and Kalai combined that extension
with the Margulis--Russo probability-derivative identity to show that
every increasing family invariant under a transitive coordinate group
has a narrow transition~\cite[Section~2]{FK96}. For graph properties,
vertex permutations act transitively on the edge coordinates, and their
theorem gives the universal width bound
$C\log(1/(2\varepsilon))/\log n$~\cite[Theorem~1.1]{FK96}.
They conjectured the denominator $(\log n)^2$ with the same dependence
on $\varepsilon$~\cite[Conjecture~1.2]{FK96}.
Theorem~\ref{thm:threshold} resolves that conjecture positively;
its numerical constant is only a convenient explicit choice.

The squared logarithm is the optimal order in $n$ when
$\varepsilon$ is fixed. Friedgut and Kalai already pointed to the
property of containing a clique with order proportional to $\log n$:
its transition width can be of order $(\log n)^{-2}$
\cite[p.~2994]{FK96}.
Bourgain and Kalai went beyond coordinate transitivity by using the
action of the symmetry group on sets of coordinates~\cite{BK97}.
For every fixed $\eta>0$ and $0<\varepsilon<1/2$, their graph-property
bound is
\[
 p_{1-\varepsilon}(\mathcal P)-p_\varepsilon(\mathcal P)
 \le C_{\eta,\varepsilon}(\log n)^{-2+\eta}.
\]
Thus their result approached the conjectured exponent arbitrarily
closely. Further progress through Fourier estimates is described below;
Friedgut gives a broader account of the influence method
in~\cite{Friedgut22}.

Two companion papers address complementary threshold-location
questions: comparison of integral and fractional covers of increasing
families~\cite[Theorem~1.1]{ExpectationEquivalence2026}, and comparison
of graph-containment thresholds with expected subgraph
counts~\cite[Theorem~1.1]{SecondKahnKalai2026}.
The present theorem controls the width of a transition for arbitrary
increasing graph properties. Its proof does not use either comparison.

Here is the stronger estimate that drives the proof. For a Boolean function $f:\{0,1\}^{E_n}\to\{0,1\}$, define its \emph{pivotal influence} at edge $e$ and its \emph{total influence} as follows. Here $x_{e\leftarrow b}$ denotes $x$ with coordinate $e$ set to $b\in\{0,1\}$, and the probability is over all coordinates other than $e$.
\[
 I_{e,p}(f)=\Pr_p\bigl(f(x_{e\leftarrow1})\ne f(x_{e\leftarrow0})\bigr),
 \qquad I_p(f)=\sum_{e\in E_n}I_{e,p}(f).
\]

\noindent\textit{Remark (notation).} $f$ is a Boolean function $f:\{0,1\}^{E_n}\to\{0,1\}$. In Theorem~\ref{thm:threshold}: $f=\1_{\mathcal P}$, where $\mathcal P$ is the nontrivial increasing graph property under consideration. In Proposition~\ref{prop:block}, Lemma~\ref{lem:block-capture}, Theorem~\ref{thm:variance}: $f$ is invariant under every permutation of $[n]$; monotonicity is not assumed. In Lemma~\ref{lem:restriction}: $f$ is an arbitrary Boolean function on $\{0,1\}^{E_n}$; neither symmetry nor monotonicity is assumed. $x$ denotes an argument configuration on the coordinate set of the function being evaluated; in $x_{e\leftarrow b}$ it is a graph configuration in $\{0,1\}^{E_n}$. In Lemma~\ref{lem:low-degree}: In the Fourier characters, $x\in\{0,1\}^{J}$; as an argument of $\Delta_i g$, $x$ ranges over the remaining coordinates $J\setminus\{i\}$. In Lemma~\ref{lem:restriction}: In $f_y(x)=f(x,y)$, $x\in\{0,1\}^{F}$ is the free-coordinate configuration, while $y\in\{0,1\}^{H}$ specifies the exterior coordinates. In the notation $x_{e\leftarrow b}$, $b\in\{0,1\}$ is the value assigned to coordinate $e$. In Lemma~\ref{lem:low-degree}: In the scalar two-point inequality, $b$ is an arbitrary real scalar. In tensorization, the same letter locally denotes the real-valued coefficient function on $\{0,1\}^{J\setminus\{i\}}$ in $g=a+b\chi_i$; this is a function of the remaining coordinates.


\begin{theorem}[Variance and influence]
\label{thm:variance}
For every integer $n\ge2$, every Boolean function $f$ on $E_n$ invariant under vertex permutations, and every $0<p<1$,
\begin{equation}
\label{eq:variance}
 \Var_p(f)\le\frac{2^{17}}{(\log n)^2}I_p(f).
\end{equation}
No monotonicity assumption is needed.
\end{theorem}

At the unbiased measure, Kelman, Kindler, Lifshitz, Minzer, and Safra
proved the graph-symmetric estimate
\[
 I_{1/2}(f)\ge
 c\,\frac{(\log n)^2}{(\log\log n)^2}\Var_{1/2}(f)
\]
for all sufficiently large $n$ and an absolute $c>0$
\cite[arXiv version~2, Section~7.4]{KKLMS20}.
This is an influence estimate at $p=1/2$; a threshold-width bound
requires control throughout the transition interval.
Theorem~\ref{thm:variance} removes the iterated-logarithm loss and
holds uniformly for every $0<p<1$, including for nonmonotone properties.

Random restrictions and symmetry of Fourier coefficients are central
to the arguments of Bourgain--Kalai and
Kelman--Kindler--Lifshitz--Minzer--Safra~\cite{BK97,KKLMS20}.
Here the restriction is chosen through the vertices: leave free all
edges touching a random block $B$ of about $\sqrt n$ vertices, while
conditioning on every other edge. Permutations inside $B$ survive
every such conditioning. Their edge orbits bound each conditional
influence, to which an elementary Fourier estimate applies.
The resulting estimate is linear in total influence, so it can be
averaged over the conditioned coordinates. Section~\ref{sec:fourier}
proves this estimate directly, including its dependence on the bias.

\noindent\textit{Remark (notation).} The notation $B$ is local to its context, denoting either the auxiliary function $T_\rho b$ or an $m$-element vertex block. In Lemma~\ref{lem:low-degree}: In the tensorization proof, $B=T_\rho b$ is a real function on the coordinates other than $i$, where $g=a+b\chi_i$. In Proposition~\ref{prop:block}: $B\in\binom{[n]}m$, with integer $3\le m\le n$. In Lemma~\ref{lem:block-capture}: The proof uses $B\in\binom{[n]}m$ with $m=\lfloor\sqrt n\rfloor$.


The second ingredient regards each edge set indexing a Fourier
coefficient as a graph. Any nonempty such graph with at most $k^2/2$
edges has a nonisolated vertex of degree at most $k$. With sufficient
probability, $B$ meets its nonisolated vertices only at that vertex.
The restriction then leaves between one and $k$ edges of this graph
free. This turns control of conditional Fourier degrees up to $k$
into control of original degrees up to $k^2/2$. Taking $k$ proportional
to $\sqrt{p(1-p)}\log n$ yields \eqref{eq:variance}; the bias factor
cancels in the estimate for the remaining Fourier degrees.
Section~\ref{sec:threshold} then integrates the probability derivative
to obtain the full quantile bound.

\noindent\textit{Remark (notation).} $k$ denotes a positive real degree cutoff, with comparisons to integer cardinalities interpreted literally rather than requiring $k$ to be an integer. In Lemma~\ref{lem:low-degree}: $k=\sigma\log m/64$, where $m\ge2$ is the real influence comparison parameter. In Lemma~\ref{lem:restriction}: $k$ is any positive real number, with no required relation to an influence comparison parameter. In Proposition~\ref{prop:block}: $k=\sigma\log m/64$, where $m$ is the integer block size satisfying $3\le m\le n$. In Lemma~\ref{lem:block-capture}, Theorem~\ref{thm:variance}: When $t=\log n\ge16$, $m=\lfloor\sqrt n\rfloor$ and $k=\sigma\log m/64$; the original Fourier levels are split at $k^2/2$.

